\chapter{Pruebas y Validación del Sistema}
\label{cap:06_pruebas_validacion}

\section{Plan de pruebas}
\pendiente{Desarrollar este apartado con metodología, resultados, decisiones y evidencias verificables.}

\section{Cobertura de pruebas}
\pendiente{Desarrollar este apartado con metodología, resultados, decisiones y evidencias verificables.}

\section{Matriz de trazabilidad entre requisitos y pruebas}
\pendiente{Desarrollar este apartado con metodología, resultados, decisiones y evidencias verificables.}

\begin{table}[htbp]\centering
\caption{Matriz de trazabilidad (ejemplo editable)}\label{tab:trazabilidad}
\begin{tabularx}{\textwidth}{@{}lXll@{}}\toprule
Requisito & Caso de prueba & Resultado & Evidencia \\\midrule
HU-XX & TC-XX: descripción & Pendiente & EV-XX \\\bottomrule
\end{tabularx}\end{table}

\section{Resultados de las pruebas}

\begin{table}[htbp]
    \centering

    \caption{Matriz de verificación de la integración del componente inteligente}
    \label{tab:resultados_integracion_ia}

    \renewcommand{\arraystretch}{1.2}

    \begin{tabularx}{\textwidth}{
        |>{\raggedright\arraybackslash}p{3cm}|
         >{\raggedright\arraybackslash}X|
         >{\raggedright\arraybackslash}X|
    }
        \hline

        \rowcolor{gray!15}
        \textbf{Componente} &
        \textbf{Comportamiento esperado} &
        \textbf{Evidencia por vincular} \\

        \hline

        \textbf{Configuración de complejidad} &
        Selección del nivel de complejidad requerido
        desde la interfaz. &
        Panel principal y solicitudes de generación. \\

        \hline

        \textbf{Generación procedural} &
        Construcción de escenarios candidatos a
        partir de los parámetros establecidos. &
        Mapas generados y pruebas funcionales
        de los algoritmos. \\

        \hline

        \textbf{Extracción de características} &
        Obtención de representaciones numéricas
        de los mapas candidatos. &
        Código fuente y pruebas del procesamiento
        de características. \\

        \hline

        \textbf{Clasificación ML} &
        Estimación del nivel de complejidad de
        los escenarios generados. &
        Resultados experimentales y pruebas de
        clasificación. \\

        \hline

        \textbf{Selección adaptativa} &
        Selección de un escenario considerando
        la complejidad solicitada y los resultados
        de evaluación. &
        Respuestas del servicio de generación
        adaptativa. \\

        \hline

        \textbf{Evaluación estructural} &
        Obtención de métricas asociadas a las
        propiedades del mapa. &
        Resultados de evaluación y pruebas
        de cálculo de métricas. \\

        \hline

        \textbf{Integración de servicios} &
        Intercambio de información entre el
        frontend y el backend. &
        Pruebas de integración y respuestas
        de la API REST. \\

        \hline

        \textbf{Presentación de resultados} &
        Visualización del escenario y de la
        información proporcionada por el backend. &
        Capturas del prototipo y pruebas
        funcionales de la interfaz. \\

        \hline

    \end{tabularx}

    \par\smallskip
    \raggedright

    \textit{Nota.} Matriz de planificación de la verificación. Los resultados y evidencias específicas deben incorporarse a partir de la ejecución registrada.

\end{table}
\FloatBarrier
\pendiente{Desarrollar este apartado con metodología, resultados, decisiones y evidencias verificables.}

\section{Pruebas con usuarios (si aplica)}
\pendiente{Desarrollar este apartado con metodología, resultados, decisiones y evidencias verificables.}
